% ═══════════════════════════════════════════════════════════════════════════════
%  CAPÍTULO — Singleton
% ═══════════════════════════════════════════════════════════════════════════════

\chapter{Singleton}

% ─── Situación Problema ──────────────────────────────────────────────────────
\section{Situación Problema}

La plataforma de Airbnb maneja una serie de parámetros de configuración
global que deben ser coherentes en toda la aplicación: la comisión de
servicio vigente, las monedas aceptadas, los umbrales de penalización por
cancelación y las políticas de reembolso activas. Estos valores los
gestiona el equipo de operaciones y pueden cambiar con el tiempo, pero en
todo momento debe existir \emph{una única versión de la verdad}.

El problema surge cuando múltiples módulos —el motor de reservas, el
servicio de pagos y el panel de anfitriones— instancian de forma
independiente su propia copia del objeto de configuración. Si en algún
momento se actualiza la comisión de servicio, solo alguna de esas copias
recibe el nuevo valor; las demás trabajan con datos desactualizados,
produciendo cobros inconsistentes y discrepancias visibles al usuario.
Una solución basada en variables globales ordinarias tampoco es suficiente,
pues no garantiza que la inicialización ocurra una única vez ni ofrece
ningún mecanismo de control de acceso concurrente.

% ─── Solución ────────────────────────────────────────────────────────────────
\section{Solución}

El patrón Singleton garantiza que la clase \texttt{ConfiguracionGlobal}
tenga exactamente una instancia durante todo el ciclo de vida de la
aplicación, y proporciona un punto de acceso único a ella. La primera
vez que cualquier módulo solicita la configuración, la instancia se crea
y se almacena internamente; las solicitudes posteriores reciben siempre
la misma referencia.

De esta manera, cuando el equipo de operaciones actualiza la comisión de
servicio, el cambio queda reflejado de inmediato para todos los módulos
sin necesidad de sincronización adicional. El motor de reservas, el
servicio de pagos y el panel de anfitriones consultan el mismo objeto,
eliminando las inconsistencias. El patrón también facilita aplicar un
mecanismo de control de acceso concurrente en entornos multihilo, ya que
la lógica de creación está centralizada en un único lugar.

% ─── Diagrama de clases ─────────────────────────────────────────────────────
\begin{figure}[H]
    \centering
    % \includegraphics[width=\textwidth]{imagenes/abstract_factory_diagrama.png}
    \caption{Diagrama de clases del patrón Abstract Factory}
\end{figure}


% ─── Roles de las Clases ────────────────────────────────────────────────────
\section{Roles de las Clases}

% Explicar la responsabilidad de cada clase/interfaz

\begin{itemize}
    \item \textbf{AbstractFactory:}
    \item \textbf{ConcreteFactory:}
    \item \textbf{AbstractProduct:}
    \item \textbf{ConcreteProduct:}
    \item \textbf{Client:}
\end{itemize}


% ─── Main ───────────────────────────────────────────────────────────────────
\section{Main}

% Explicar qué realiza el programa principal

\begin{lstlisting}[language=Java, caption=NombreClase]
 
\end{lstlisting}


% ─── Salida del Main ────────────────────────────────────────────────────────
\section{Salida del Main}

\begin{lstlisting}[language=Java, caption=NombreClase]
 
\end{lstlisting}


% ─── Clases ─────────────────────────────────────────────────────────────────
\section{Clases}

% ─── Clase 1 ────────────────────────────────────────────────────────────────
\subsection{NombreClase}

\begin{lstlisting}[language=Java, caption=NombreClase]
 
\end{lstlisting}


% ─── Clase 2 ────────────────────────────────────────────────────────────────
\subsection{NombreClase}

\begin{lstlisting}[language=Java, caption=NombreClase]
 
\end{lstlisting}


% ─── Clase 3 ────────────────────────────────────────────────────────────────
\subsection{NombreClase}

\begin{lstlisting}[language=Java, caption=NombreClase]
 
\end{lstlisting}


% ─── Conclusiones ───────────────────────────────────────────────────────────
\section{Conclusiones}

% Explicar:
% - Ventajas del patrón
% - Cuándo usarlo
% - Beneficios obtenidos
% - Posibles desventajas
% - Relación con principios SOLID